\documentclass[a4paper,addpoints]{article}
\usepackage{amsfonts,amssymb,amsmath, amsthm}
\usepackage{xcolor}
\usepackage{pgf,tikz,pgfplots}
\pgfplotsset{compat=1.15}
\usepgfplotslibrary{fillbetween}
\usepackage{mathrsfs}
\usetikzlibrary{arrows}
\usetikzlibrary{calc}
\usepackage{tabu}

\usepackage{tikz}
\usetikzlibrary {shapes , arrows , fit , calc, positioning}
\tikzset{box/.style={draw, rectangle, rounded corners, thick, node distance =7 em, text width =6 em, text centered, minimum height =2.5 em}}
\tikzset{line/.style={draw, thick, -latex'}}

\usepackage{tikz-uml}
\tikzumlset{actor below=0.4cm, fill usecase=white, fill class=white, text=black, draw=black}


\begin{document}

Il diagramma è composto da 3 principali personaggi: \textbf{un utente non registrato}, \textbf{un utente registrato} e \textbf{il GPS}, che fornisce la posizione dell'utente e delle offerte. L'utente non registrato può solo registrarsi o fare il log in; non inserisce alcun dato finché non è loggato. Le funzionalità principali, riservate all'utente registrato, sono la visualizzazione delle offerte (che include la visualizzazione della posizione), l'ordinamento per prezzo, la ricerca dinamica e il filtro per marca, per poter vedere ciò che interessa davvero.

\begin{center}
\begin{tikzpicture}
    \begin{umlsystem}[x=4] {Sistema}
        \umlusecase[x=4, y=-1.5, width=1.8cm]{visualizzazione di offerte} 
        \umlusecase[x=4, y=-7, width=1.8cm]{visualizzazione della posizione} 
        \umlusecase[x=4, width=1.8cm] {registrazione e log in}
        \umlusecase[x=2, y =-5.5, width=1.8cm] {ordinare in base a prezzo}       
        \umlusecase[x=-1, y =-6, width=1.8cm] {ricerca dinamica}    
        \umlusecase[x=0, y=-1.5, width=1.8cm]{visualizzazione prezzi nel tempo} 
        \umlusecase[x=7, y=-4.5, width=1.8cm]{filtrare per marca} 
    \end{umlsystem}
        
    \umlactor{cliente non registrato} 
    \umlactor[y=-4] {cliente registrato}
    \umlactor[x=12, y=-7] {gps}

    \umlinherit{cliente registrato}{cliente non registrato}

    % l'offerta include la posizione (GPS)
    \umlinclude{usecase-1}{usecase-2}
    \umlinclude{usecase-1}{usecase-5}
    \umlinclude{usecase-1}{usecase-4}
    \umlinclude{usecase-1}{usecase-7}

    \umlassoc{cliente non registrato}{usecase-3}
    
    \umlassoc{cliente registrato}{usecase-1}
    \umlassoc{cliente registrato}{usecase-4}
    \umlassoc{cliente registrato}{usecase-5}
    \umlassoc{cliente registrato}{usecase-7}
    \umlassoc{cliente non registrato}{usecase-6}

    \umlassoc{gps}{usecase-2}
\end{tikzpicture}
\end{center}

\end{document} 
